\documentclass[12pt,a4paper]{article}
\usepackage[utf8]{inputenc}
\usepackage[T1]{fontenc}
\usepackage[english]{babel}
\usepackage{lmodern,geometry,setspace,amsmath,amssymb,booktabs,array,longtable}
\usepackage{fancyhdr,listings,xcolor}
\geometry{top=2cm,bottom=2cm,left=2.5cm,right=2.5cm}
\setstretch{1.2}
\pagestyle{fancy}\fancyhf{}\rfoot{\thepage}
\definecolor{codebg}{RGB}{246,247,249}
\lstset{language=Python,basicstyle=\ttfamily\footnotesize,backgroundcolor=\color{codebg},
  frame=single,rulecolor=\color{gray!40},breaklines=true,showstringspaces=false,
  columns=fullflexible,keepspaces=true,
  literate={Ă}{{\u{A}}}1 {Â}{{\^A}}1 {Î}{{\^I}}1 {Ș}{{\c{S}}}1 {Ț}{{\c{T}}}1
           {ă}{{\u{a}}}1 {â}{{\^a}}1 {î}{{\^\i}}1 {ș}{{\c{s}}}1 {ț}{{\c{t}}}1}
\newcommand{\file}[1]{\texttt{#1}}
\renewcommand{\contentsname}{Cuprins}

\begin{document}
\begin{titlepage}
\begin{center}
\textbf{MINISTERUL EDUCAȚIEI ȘI CERCETĂRII AL REPUBLICII MOLDOVA}\\[0.5em]
\textbf{UNIVERSITATEA TEHNICĂ A MOLDOVEI}\\
Facultatea Calculatoare, Informatică și Microelectronică\\
Departamentul Ingineria Software și Automatică
\vspace{3cm}

{\Large\bfseries RAPORT}\\[0.7em]
{\large Lucrare de laborator nr. 3}\\[0.5em]
{\large Cifruri polialfabetice. Algoritmul Playfair}\\[0.7em]
Disciplina: Criptografie\\
Varianta 23 --- sarcina 3.1
\vspace{2.5cm}

\begin{flushright}
Realizat de: Tacu Igor, FAF-242\\[0.5em]
Lector: Maia Zaica
\end{flushright}
\vfill
\textbf{Chișinău --- 2026}
\end{center}
\end{titlepage}

\pagenumbering{roman}
\tableofcontents
\newpage
\pagenumbering{arabic}

\section{Sarcina și alegerea variantei}

Numărul meu în registru este \textbf{23}. Conform notei de la sfârșitul îndrumarului,
numerele impare rezolvă \textbf{sarcina 3.1}: implementarea cifrului Playfair pentru
mesaje în limba română. Programul trebuie să primească operația, cheia și mesajul
sau criptograma, apoi să afișeze rezultatul. Cheia are minimum șapte litere;
caracterele nepermise sunt respinse cu o indicație despre intervalul admis.
Spațiile din textul decriptat se restabilesc manual. Sarcina 3.2, despre Vigenère,
revine numerelor pare și nu face parte din această variantă.

\section{Modelul folosit}

\subsection{Alfabetul și matricea}

Am păstrat ordinea alfabetului românesc folosită în laboratorul 1:
\begin{center}
\small\texttt{A Ă Â B C D E F G H I Î J K L M N O P Q R S Ș T Ț U V W X Y Z}
\end{center}
Sunt 31 de litere. Enunțul menționează simultan 31 de litere românești și intervalul
\texttt{A-Z/a-z}; ultimul, luat literal, ar conține numai 26 de litere. Pentru a putea
introduce mesaje românești, programul acceptă și cele cinci litere cu diacritice,
inclusiv variantele cu sedilă pentru Ș și Ț. Acceptă spații în mesajul clar,
deoarece etapa de pregătire din îndrumar le elimină. Cheia și criptograma nu
acceptă spații. Cifrele, punctuația, taburile și alte simboluri sunt respinse.

O matrice dreptunghiulară completă cu 31 de poziții ar avea doar dimensiunile
$1\times31$ sau $31\times1$. Acestea nu ar permite folosirea efectivă a regulii
dreptunghiului. Îndrumarul admite explicit o matrice cu \emph{aproape toate}
literele și restabilirea intuitivă a literelor eliminate. Am ales \textbf{5 rânduri
și 6 coloane}: 30 de poziții. Litera \texttt{Q}, rară în cuvintele românești,
este asimilată cu \texttt{K} înainte de cifrare. Este o pierdere de informație:
după descifrare, un \texttt{K} provenit din \texttt{Q} poate fi recunoscut numai din
context. Nu elimin diacriticele, pentru că ele sunt esențiale pentru un text
românesc lizibil.

\subsection{Pregătirea textului}

Mai întâi transform literele în majuscule, elimin spațiile, normalizez formele
Unicode echivalente și aplic \texttt{Q}$\to$\texttt{K}. Apoi formez perechi de câte
două litere. Dacă următoarele două litere sunt identice, inserez \texttt{X} între
ele; pentru perechea \texttt{XX} folosesc \texttt{Z}, ca să nu produc din nou o
pereche egală. Aceeași regulă completează mesajele cu lungime impară. De exemplu:
\[
\texttt{ANA ARE MERE}\longrightarrow\texttt{AN AX AR EM ER EX}.
\]
Regula se aplică perechilor formate efectiv, nu oricăror două litere identice
care devin vecine după eliminarea spațiilor. Literele de umplere nu se șterg
automat la decriptare: un \texttt{X} sau \texttt{Z} poate aparține mesajului
original, iar programul nu poate decide sigur acest lucru.

\subsection{Construirea cheii și regulile Playfair}

Cheia de demonstrație este \texttt{SECURITATE} (10 litere). Repetițiile se elimină
după prima apariție: \texttt{SECURITA}. Adaug apoi, în ordinea de mai sus, toate
literele alfabetului de 30 care încă lipsesc. Matricea rezultată este:

\begin{center}
\renewcommand{\arraystretch}{1.3}
\begin{tabular}{cccccc}
\toprule
S&E&C&U&R&I\\
T&A&Ă&Â&B&D\\
F&G&H&Î&J&K\\
L&M&N&O&P&Ș\\
Ț&V&W&X&Y&Z\\
\bottomrule
\end{tabular}
\end{center}

Pentru o digramă, pozițiile sunt $(r_1,c_1)$ și $(r_2,c_2)$:
\begin{itemize}
\item același rând: ambele coloane avansează cu o poziție modulo 6;
\item aceeași coloană: ambele rânduri avansează cu o poziție modulo 5;
\item rânduri și coloane diferite: rezultatul este $(r_1,c_2)$, respectiv
  $(r_2,c_1)$, adică celelalte colțuri ale dreptunghiului.
\end{itemize}
La decriptare, primele două deplasări merg în sens opus; regula dreptunghiului
este aceeași. Exemple simple din matricea de mai sus: \texttt{SE}$\to$\texttt{EC}
(rând), \texttt{ST}$\to$\texttt{TF} (coloană), \texttt{SA}$\to$\texttt{ET}
(dreptunghi). Trecerea peste margine se vede în \texttt{IS}$\to$\texttt{SE}.

\section{Implementare}

Programul este în \file{lab3/playfair.py}, Python 3, fără pachete externe.
Funcțiile sunt separate după responsabilitate: \file{clean} validează și
normalizează intrarea; \file{key\_square} construiește matricea;
\file{prepare} produce digramele; \file{transform} aplică regulile;
\file{encrypt} și \file{decrypt} sunt operațiile folosite de interfața de consolă.

Partea decisivă a transformării este scurtă și oglindește exact cele trei cazuri:
\begin{lstlisting}
if ra == rb:
    output.extend((square[ra * COLS + (ca + direction) % COLS],
                   square[rb * COLS + (cb + direction) % COLS]))
elif ca == cb:
    output.extend((square[((ra + direction) % ROWS) * COLS + ca],
                   square[((rb + direction) % ROWS) * COLS + cb]))
else:
    output.extend((square[ra * COLS + cb], square[rb * COLS + ca]))
\end{lstlisting}
Valoarea \file{direction} este $+1$ la criptare și $-1$ la decriptare. Operatorul
\file{\%} face rotația la capătul rândului sau al coloanei. Pozițiile literelor
sunt construite o singură dată într-un dicționar, înainte de parcurgerea
digramelor; astfel nu este necesară căutarea repetată în matrice.

La criptare se validează cheia și mesajul, se afișează matricea și digramele
pregătite, apoi criptograma. La decriptare se verifică în plus lungimea pară a
criptogramei și absența perechilor identice, imposibile într-o criptogramă
Playfair validă. Programul afișează textul pregătit rezultat și amintește că
spațiile, eventualele umpleri și litera \texttt{Q} trebuie restabilite manual.

\section{Exemplu complet de criptare și decriptare}

Folosesc cheia \texttt{SECURITATE} și mesajul \texttt{ȘTIINȚĂ ȘI CERCETARE}.
După normalizare și formarea digramelor obțin:
\begin{center}
\texttt{ȘT IX IN ȚĂ ȘI CE RC ET AR EX}
\end{center}
\texttt{X} a fost inserat între cele două litere \texttt{I} consecutive,
iar ultimul \texttt{E} a fost completat cu \texttt{X}. Următorul tabel arată
toate substituțiile, astfel încât rezultatul poate fi verificat manual:

\begin{center}
\begin{tabular}{ccccccccccc}
\toprule
Clar&ȘT&IX&IN&ȚĂ&ȘI&CE&RC&ET&AR&EX\\
Cifrat&LD&UZ&CȘ&WT&ZD&UC&IU&SA&BE&UV\\
\bottomrule
\end{tabular}
\end{center}

Criptograma fără separatori este:
\begin{center}
\large\texttt{LDUZCȘWTZDUCIUSABEUV}
\end{center}
Decriptarea cu aceeași cheie dă \texttt{ȘTIXINȚĂȘICERCETAREX}. După eliminarea
celor două litere de umplere și plasarea manuală a spațiilor, recuperez
\texttt{ȘTIINȚĂ ȘI CERCETARE}. Exemplul arată de ce rezultatul brut al
decriptării nu este întotdeauna identic cu forma inițială a mesajului.

\section{Verificări și cazuri de eroare}

Am rulat \file{python3 -m unittest -v test\_playfair.py}. Cele șapte teste acoperă
matricea (30 de litere distincte), cele trei reguli și rotația la margini,
pregătirea literelor duble, mesaje impare, diacriticele, validarea datelor și
relația
\[
 D_k(E_k(m))=\operatorname{prepare}(m).
\]
Unul dintre teste verifică separat toate cele $30\cdot29=870$ de digrame cu
litere diferite: cifrează fiecare digramă și confirmă că decriptarea o recuperează.
Egalitatea este cu textul \emph{pregătit}, deoarece inserările și pierderea lui
\texttt{Q} nu pot fi inversate automat. Am verificat această proprietate pentru
mesaje cu diacritice, cu litere \texttt{X} repetate și cu \texttt{Q}.

\begin{center}
\begin{tabular}{p{0.36\textwidth}p{0.49\textwidth}}
\toprule
Intrare sau situație&Rezultat urmărit\\
\midrule
Cheia \texttt{SCURT}&Respinsă: are sub 7 litere.\\
Mesajul \texttt{ABC!}&Respins: semnul \texttt{!} nu este permis; se afișează caracterele admise.\\
Criptograma \texttt{ABC}&Respinsă: lungime impară.\\
Criptograma \texttt{AA}&Respinsă: digramă cu litere identice.\\
Textul \texttt{ştiinţă}&Normalizat la \texttt{ȘTIINȚĂ}.\\
\bottomrule
\end{tabular}
\end{center}

\section{Limite și concluzii}

Playfair ascunde frecvența literelor individuale mai bine decât un cifru de
substituție monoalfabetic, deoarece criptează \emph{perechi} și o literă poate
avea rezultate diferite în funcție de vecina sa. Totuși, perechile repetate
dintr-un text lung lasă tipare, iar algoritmul este nepotrivit pentru protecția
datelor actuale. Implementarea de aici are scop didactic.

Varianta de 30 de poziții respectă recomandarea din îndrumar pentru un alfabet
care nu încape într-o matrice potrivită, dar \texttt{Q}$\to$\texttt{K} înseamnă
că descifrarea nu poate reconstitui automat toate mesajele. Același lucru se
aplică spațiilor și literelor de umplere. Programul lasă aceste decizii la
cititor, în acord cu cerința de restabilire manuală a formei finale.

\section*{Fișiere și sursă}
\addcontentsline{toc}{section}{Fișiere și sursă}
\begin{itemize}
\item \file{lab3/playfair.py} --- implementarea și interfața de consolă;
\item \file{lab3/test\_playfair.py} --- verificările automate;
\item \file{lab3/report.tex} --- sursa acestui raport.
\end{itemize}
Îndrumar: \emph{Lucrare de laborator nr. 3. Cifruri polialfabetice},
secțiunea 3.2 și sarcina 3.1, paginile 3--7 (documentul furnizat pentru curs).

\end{document}
